% Using InitechBase, Appendix E: Data Types in Expressions (PUB-REF-0001
% CP-11: generated from spec/samir/xbase_coercion.json by
% harness/pubs/samir_tables.c coercion; the row notes are defined here, one
% per rule name of the JSON, and an undefined one stops the build).
\chapter{Data Types in Expressions}
\Gates{test-samir-spec, test-xbase-eval, test-xbase-coercion}
\NoDictionaryHeads
The operators of InitechBase work on values of particular types.  The first
table gives, for each operator, the result of each combination of types that
it accepts, and the message for some combinations it refuses.  A combination
not listed is refused as well, with \xmsg{9}{Data type mismatch.}  The second
table gives the result of storing a value of each type in a field of each
type with \cmd{REPLACE}.

\DefineCoercionRule{R_concat_plus}{Joined as they are}
\DefineCoercionRule{R_concat_minus}{Trailing blanks of the left string moved to the end}
\DefineCoercionRule{R_date_plus_num}{Days added or subtracted; a blank date stays blank}
\DefineCoercionRule{R_date_minus_date}{Number of days; 0 when a date is blank}
\DefineCoercionRule{R_blankdate_high}{A blank date is later than any date}
\DefineCoercionRule{R_order_same_type}{Both sides must be of one type}
\DefineCoercionRule{R_substr}{True when the left string occurs in the right}
\DefineCoercionRule{R_no_truthiness}{A number is not a logical value}
\DefineCoercionRule{R_begins_with}{SET EXACT OFF: the left begins with the right}
\DefineCoercionRule{R_exact_blankpad}{SET EXACT ON: equal but for trailing blanks}
\DefineCoercionRule{R_same_type}{Both sides must be of one type}

\section{Operators}
\begingroup\small
\input{samir_coercion}
\endgroup

\draft{the row Memo / Character: this version does not store a character
value replaced into a memo field (see REPLACE in Chapter~5)}

\section{Comparing Character Strings}
For \texttt{=}, \texttt{<>} and \texttt{\#} between character strings the
result depends on SET EXACT.  With SET EXACT OFF, the comparison is true when
the left string begins with the right one, so an empty string on the right
equals every string.  With SET EXACT ON, the strings must be equal, apart from
trailing blanks.  \texttt{<}, \texttt{>}, \texttt{<=} and \texttt{>=} compare the
codes of the characters one by one, whatever the setting.
